\documentclass[10pt]{beamer}
%================================================
%---------------PACOTES
%================================================
\usepackage{appendixnumberbeamer}
\usepackage{booktabs}
\usepackage[scale=2]{ccicons}
\usepackage{pgfplots}
\usepackage{xspace}
\newcommand{\themename}{\textbf{\textsc{metropolis}}\xspace}
\renewcommand{\raggedright}{\leftskip=0pt \rightskip=0pt plus 0cm}
\usepackage[brazil]{babel}
\usepackage[T1]{fontenc}
\usepackage[utf8]{inputenc}
\usepackage{fixltx2e}
\usepackage{graphicx}
\usepackage{pdflscape}
\usepackage{epstopdf}
\usepackage{pdfpages}
\usepackage{verbatim}
\tolerance=1000
\usepackage[outputdir=build/]{minted}
\usepackage{xcolor}
\usepackage[all]{xy}
\usepackage{times}
\usepackage{tikz}
\usepackage{ragged2e}
\usepackage{caption}
\usepackage{subcaption}
\usepackage{float}

\usepackage{indentfirst}
\usepackage{graphicx}
\usepackage[timeinterval=1]{tdclock}

\usepackage{multirow}
\usepackage{multicol}
\usepackage{lipsum}
\usepackage{ragged2e}
\newcommand{\jus}{\justifying}
\apptocmd{\frame}{}{\justifying}{}
\apptocmd{\block}{}{\justifying}{}
\apptocmd{\column}{}{\justifying}{}
\usepackage{enumerate}
\usepackage{handoutWithNotes}
%\pgfpagesuselayout{1 on 1 with notes landscape}[a4paper,border shrink=5mm]
\hypersetup{colorlinks,citecolor=verde,linkcolor=orange,urlcolor=verde}
\usepackage[alf]{abntex2cite}%Citações padrão ABNT}
%=============================
\usetheme[sectionpage=none, subsectionpage=progressbar]{metropolis}
\usepgfplotslibrary{dateplot}
\useoutertheme[subsection=false,footline=authortitle]{miniframes}
\setbeamertemplate{footline}[frame number]%Numerando as Páginas
\hypersetup{pdfpagemode=FullScreen}%Visualizar em Tela cheia
\hypersetup{pdfpagelayout=SinglePage}%Layout da pagina
%-----Definição de novos comandos
\providecommand{\sin}{} \renewcommand{\sin}{\hspace{2pt}\textrm{sen}}
\providecommand{\tan}{} \renewcommand{\tan}{\hspace{2pt}\textrm{tg}}
\setbeamerfont{section in toc}{size=\small}
\setbeamerfont{subsection in toc}{size=\scriptsize}%\footnotesize
\setbeamerfont{subsubsection in toc}{size=\scriptsize}
\AtBeginEnvironment{minted}{\renewcommand{\fcolorbox}[4][]{#4}}
\AtBeginSection[]{\begin{frame}<beamer>\frametitle{Seção atual}\tableofcontents[currentsection]\end{frame}}
\LetLtxMacro{\OldIncludegraphics}{\includegraphics}
\renewcommand{\includegraphics}[2][]{\OldIncludegraphics[height=0.8\textheight,keepaspectratio, #1]{#2}}
\newcommand{\dis}{\displaystyle}
\newcommand{\lb}{\left\lbrace}
\newcommand{\rb}{\right\rbrace}
%================================================
%---------------CORES DO TEMPLATE
%================================================
\definecolor{verde}{HTML}{008069}
\definecolor{verdepastel}{HTML}{d8f3dc}
\definecolor{verdesolution}{HTML}{beff9e}
\definecolor{vsplbggray}{HTML}{f4f8ff}
%-----------------------------------------------
\setbeamercolor{alerted text}{fg=orange}
\setbeamercolor*{palette primary}{fg=verde!60!black,bg=gray!30!white}
\setbeamercolor*{palette tertiary}{bg=verde!90,fg=white}
\setbeamercolor*{palette quaternary}{fg=verde,bg=gray!5!white}
\setbeamercolor*{sidebar}{fg=black,bg=black}
\setbeamercolor{progress bar}{fg=verde}%Linha q passa as seções
\setbeamercolor*{titlelike}{parent=palette primary}
\setbeamercolor{titlelike}{parent=palette primary,fg=verde}%-Modifica a cor do título
\setbeamercolor{frametitle}{bg=verdepastel,fg=verde}
\setbeamercolor{frametitle right}{bg=black}
\setbeamercolor{background canvas}{bg=white}
%================================================
%---------------BLOCOS  
%================================================
\setbeamertemplate{blocks}[rounded][shadow]
\uselanguage{Brazil}
\languagepath{Brazil}
\deftranslation[to=Brazil]{Theorem}{Teorema}
\deftranslation[to=Brazil]{Problem}{Problema:}
\deftranslation[to=Brazil]{Example}{Solução:}
\deftranslation[to=Brazil]{Solution}{Solução:}
\setbeamercolor{block title}{fg=vsplbggray,bg=verde}
\setbeamercolor{block body}{bg=white}
\setbeamercolor{block title alerted}{bg=alerted text.fg!20}
\setbeamercolor{block body alerted}{bg=alerted text.fg!10}
%----------Problema--------------
\definecolor{vblock}{HTML}{E30022}
\BeforeBeginEnvironment{problem}{
\setbeamercolor{block title}{fg=vblock,bg=red!30!white}
\setbeamercolor{block body}{fg=vblock, bg=red!15!white}
}
\AfterEndEnvironment{problem}{
 \setbeamercolor{block title}{use=structure,fg=structure.fg,bg=structure.fg!20!bg}
 \setbeamercolor{block body}{parent=normal text,use=block title,bg=block title.bg!50!bg, fg=black}}
%----------Solução--------------
\BeforeBeginEnvironment{solution}{%
\setbeamercolor{block title}{fg=verde,bg=verdesolution}
\setbeamercolor{block body}{fg=verde, bg=green!5!white}
}
\AfterEndEnvironment{solution}{
 \setbeamercolor{block title}{use=structure,fg=structure.fg,bg=structure.fg!20!bg}
 \setbeamercolor{block body}{parent=normal text,use=block title,bg=block title.bg!50!bg, fg=black}
}
%----------Exemplo--------------
\setbeamercolor{block title example}{use=example text,fg=example text.fg,bg=example text.fg!50!bg}
\setbeamercolor{block body example}{parent=normal text,use=block title example,bg=block title example.bg!50!bg}
%================================================
%---------------CAPA
%================================================
\title{\bfseries Business Model Innovation for Sustainability:\\ A Value Change for Survival}
\subtitle{\textcolor{orange}{Master Thesis' Defense}}
\date{\textcolor{black}{\today}}
\date{\textcolor{brown}{28th July 2022}}
\author{\bf
Student: Pietro Mecca\\
Supervisor: Luigi Nasta  \\
Co-supervisor: Luca Pirolo 
}
 \institute{\textcolor{verde}{Master in Management at Libera Università Degli Studi Sociali Guido Carli}}
\titlegraphic{
\begin{picture}(0,0)
\put(305,-115){\makebox(0,0)[rt]{\includegraphics[width=3.5cm]{lulo.jpg}}}
\end{picture}
}
\begin{document}
\maketitle
%================================================
%---------------Sumário
%================================================
\begin{frame}{Disposition}
\hypersetup{linkcolor=black}
\setbeamertemplate{section in toc}[sections numbered]
\tableofcontents
\end{frame}
%================================================
%---------------Elementos Textuais
%================================================
\section{Purpose}

\subsection{O problema}
\begin{frame}{Problema}
\lipsum[1-1]
\end{frame}


\subsection{Motivação}
\begin{frame}{Blocos e citações}
\begin{block}{Como citar?}
\begin{table}[H]
\begin{tabular}{|l|l|l|}
\hline
Tipo de citação& Comando &  Resultado    \\ \hline
Citação ativa &\ citeonline{lamport1994latex}&  \citeonline{lamport1994latex}   \\ \hline
Citação passiva &cite{lamport1994latex}  &\cite{lamport1994latex}     \\ \hline
\end{tabular}
\end{table}
\end{block}
%---------Problemas e soluções
\begin{problem}
Você pode usar este bloco para abordar determinado problema exitente.
\end{problem}
\begin{solution}
Ideia, solução.
\end{solution}    
\end{frame}

\subsection{Objetivos}
\begin{frame}{Objetivos}
Escrever no infinitivo.
\begin{itemize}
\item Estender ; 
\item Fazer ;
\item Mostrar.
\end{itemize}
\begin{block}{Objetivos}
Os objetivos orientarão a apresentação não só da metodologia, mas, sobretudo, dos resultados, da discussão e conclusão, devem ser definidos de forma que facilite a leitura e permita uma avaliação fácil e rápida.
\end{block}
\end{frame}

\begin{frame}{New}
\lipsum[1-1]
\end{frame}

\section{Revisão da literatura}
\section{Tópicos de pesquisa}
\section{Considerações finais}
\begin{frame}{New}
\lipsum[1-1]
\end{frame}
%================================================
%---------------Elementos Pós-textuais
%================================================
{\setbeamercolor{palette primary}{fg=black, bg=white}
\begin{frame}[standout]
 \begin{figure}[H]
\centering
\includegraphics[width=1.5\linewidth]{images/gratidao.png}
\end{figure}
\end{frame}
}  
%================================================
%----------------BIBLIOGRAFIA
%================================================
\section*{Bibliografia}
\begin{frame}[allowframebreaks]{Referências bibliográficas}
%\nocite{*}
\bibliography{bibliografia}
\end{frame}
%================================================
%----------------FIM DO TEMPLATE
%================================================
\end{document}
